\section{Introduction}\label{sec:intro}

Open-source game development represents a distinct engineering domain characterized by multidisciplinary assets (game logic, graphics, audio, shaders, and level data) and predominantly volunteer-driven development cycles~\cite{marsad2026ossgamebench}. When users and contributors encounter defects, issues are tracked in repository issue trackers, investigated, and addressed through pull requests (PRs)~\cite{gousios2014pull}. While many bug reports are resolved with a single PR, a subset accumulates multiple linked PRs over time. This multi-PR pattern can arise from various development workflows: modular decomposition of fixes across subsystems, abandoned initial attempts replaced by subsequent patches, or separate fixes for newly introduced regressions.

From a project management and triage perspective, maintainers naturally seek diagnostic signals: if issues attracting multiple PRs are systematically slower to close, observing a second PR could serve as an actionable trigger for maintainer re-triage. However, interpreting this relationship requires caution due to a fundamental methodological challenge: \emph{reverse causality via time-dependent link accumulation}. Because pull requests can only be linked to an issue while that issue remains open, longer-lived issues have substantially more calendar opportunity to accumulate PR references. An observed statistical association between multi-PR issues and prolonged closure times might simply reflect the passage of time rather than higher intrinsic defect difficulty.

Furthermore, empirical repository studies must rigorously separate observable trace data from unmeasured operational concepts:
\begin{enumerate}
  \item \textbf{Calendar Duration vs.\ Debugging Effort}: The elapsed calendar duration ($\texttt{closed\_at} - \texttt{created\_at}$) measures wall-clock time, which conflates active debugging with triage queues, contributor dormancy, review latency, and volunteer availability~\cite{mockus2002two}.
  \item \textbf{Issue Closure vs.\ Verified Fix}: Issue closure indicates an administrative state transition, which does not guarantee that the defect was successfully resolved or verified against regressions; issues are also closed due to obsolescence, duplicates, or lack of reproducibility~\cite{herzig2013not}.
  \item \textbf{Reported Bugs vs.\ Actual Defect Occurrence}: Issue trackers record reported and labelled defects, omitting unfiled bugs or informal patches.
\end{enumerate}

\textbf{Research Question (RQ).} We formulate our central inquiry:
\begin{quote}
\emph{Among closed, labelled bug issues in open-source games, is multi-PR resolution associated with longer issue closure times, and does this relationship reflect genuine workflow delay or time-dependent link accumulation?}
\end{quote}

\textbf{Contribution beyond prior work.} While the original OSSGameBench publication~\cite{marsad2026ossgamebench} introduced the dataset and described high-level repository distributions, it did not investigate bug resolution dynamics. This coursework study provides:
\begin{itemize}
  \item A rigorous cohort filtering pipeline isolating 63,688 closed strict bugs across 970 repositories.
  \item A rank-based comparison of single-PR vs.\ multi-PR closure times with project-clustered uncertainty estimation.
  \item An empirical PR-arrival timing analysis demonstrating that secondary PRs accumulate predominantly at the end of issue lifespans.
  \item A fixed-landmark survival analysis ($L=30$ days) that blocks the link-accumulation channel and evaluates the residual association.
  \item A 100-issue random sample audit evaluating bug-label precision.
  \item Actionable, grounded recommendations for open-source game maintainers, with explicit empirical boundaries.
\end{itemize}
